What does it mean for an AI system to understand a request, recognize an intention, or know when it has misunderstood? These questions reach into the essence of intelligence—and acquire urgency when systems act as autonomous agents. An instruction can be followed literally while its purpose is missed. A fluent answer can conceal uncertainty. Understanding a harmful request perfectly can itself enable harm.
This keynote argues that the science of language, meaning, understanding, and interaction should fundamentally reorient NLP’s research agenda and inform how we assess AI’s capabilities and risks.
Arabic offers a particularly productive setting for developing that science. Rich morphology, unwritten vowels, dialect diversity, diglossia, script variation, and code-switching converge within Arabic communicative encounters shaped by social relationships and historical knowledge. I introduce the Arabic Conjunction Hypothesis to investigate how these dimensions jointly shape interpretation and communicative action. Arabic’s connections with Persian, Urdu, and Turkic traditions extend the inquiry across languages, scripts, and communities. Arabic’s scientific value lies in this productive conjunction, without requiring a claim of exceptionalism.
The resulting questions connect directly to debates about alignment and control: whose intentions should a system follow, what evidence supports its interpretation, and when should it ask, defer, challenge, or revise? Understanding is necessary for many forms of trustworthy assistance, but it does not establish benevolent goals or legitimate authority. A science of meaning and interaction can help make those distinctions explicit and testable.
Through this lens, annotation, evaluation, and dialogue become instruments for discovering how understanding works. The shift from NLP for Arabic to NLP through Arabic makes Arabic a source of insights that can reshape the field’s questions, methods, and foundations. What would we build—and what would we measure—if understanding in interaction became NLP’s organizing purpose? Pursuing that question also challenges how we define and evaluate intelligence in AI more broadly.
If we fear what AI might do, we must examine what it understands, whose purposes it serves, and what warrants trusting it to act.